\documentclass[11pt, a4paper]{article}

% ── Codificación y fuentes ─────────────────────────────────────────────────
\usepackage[utf8]{inputenc}
\usepackage[T1]{fontenc}
\usepackage[spanish]{babel}
\usepackage{lmodern}

% ── Márgenes y geometría ──────────────────────────────────────────────────
\usepackage[top=2.5cm, bottom=2.5cm, left=2.8cm, right=2.8cm]{geometry}

% ── Color y gráficos ──────────────────────────────────────────────────────
\usepackage[table]{xcolor}
\definecolor{grisclaro}{HTML}{F5F5F5}
\definecolor{grismedio}{HTML}{CCCCCC}
\definecolor{grisoscuro}{HTML}{444444}
\definecolor{negrotexto}{HTML}{1A1A1A}
\definecolor{rojoacento}{HTML}{C0392B}
\definecolor{verdestatus}{HTML}{1E7A42}

% ── Tipografía de secciones ───────────────────────────────────────────────
\usepackage{titlesec}
\titleformat{\section}
  {\large\bfseries\color{negrotexto}}
  {\thesection.}{0.6em}{}[\vspace{-4pt}\rule{\linewidth}{0.6pt}\vspace{2pt}]

\titleformat{\subsection}
  {\normalsize\bfseries\color{grisoscuro}}
  {\thesubsection.}{0.5em}{}

\titleformat{\subsubsection}
  {\normalsize\itshape\color{grisoscuro}}
  {\thesubsubsection.}{0.5em}{}

% ── Encabezado y pie de página ────────────────────────────────────────────
\usepackage{fancyhdr}
\pagestyle{fancy}
\fancyhf{}
\fancyhead[L]{\small\textit{ATARA — Requisitos de Despliegue}}
\fancyhead[R]{\small\textit{v1.0 · 2026}}
\fancyfoot[C]{\small\thepage}
\renewcommand{\headrulewidth}{0.4pt}

% ── Tablas ────────────────────────────────────────────────────────────────
\usepackage{booktabs}
\usepackage{tabularx}
\usepackage{array}
\newcolumntype{L}[1]{>{\raggedright\arraybackslash}p{#1}}
\newcolumntype{C}[1]{>{\centering\arraybackslash}p{#1}}

% ── Bloques de código y variables ─────────────────────────────────────────
\usepackage{listings}
\lstset{
  basicstyle=\ttfamily\small,
  backgroundcolor=\color{grisclaro},
  frame=single,
  framesep=4pt,
  rulecolor=\color{grismedio},
  breaklines=true,
  columns=fullflexible,
  keepspaces=true,
  showstringspaces=false,
}

% ── Cajas informativas ────────────────────────────────────────────────────
\usepackage{tcolorbox}
\tcbuselibrary{skins, breakable}

\newtcolorbox{notabox}{
  colback=grisclaro,
  colframe=grismedio,
  fonttitle=\bfseries\small,
  title=Nota,
  left=6pt, right=6pt, top=4pt, bottom=4pt,
  arc=3pt,
}

\newtcolorbox{warningbox}{
  colback=yellow!8!white,
  colframe=yellow!60!black,
  fonttitle=\bfseries\small,
  title=Advertencia,
  left=6pt, right=6pt, top=4pt, bottom=4pt,
  arc=3pt,
}

% ── Hiperlinks ────────────────────────────────────────────────────────────
\usepackage[hidelinks]{hyperref}
\usepackage{xurl}

% ── Listas ───────────────────────────────────────────────────────────────
\usepackage{enumitem}
\setlist[itemize]{leftmargin=1.4em, topsep=3pt, itemsep=1pt}
\setlist[enumerate]{leftmargin=1.4em, topsep=3pt, itemsep=1pt}

% ── Espaciado entre párrafos ──────────────────────────────────────────────
\setlength{\parskip}{5pt}
\setlength{\parindent}{0pt}

% ═════════════════════════════════════════════════════════════════════════
\begin{document}

% ── Portada ──────────────────────────────────────────────────────────────
\begin{titlepage}
  \centering
  \vspace*{3cm}
  {\huge\bfseries ATARA}\\[0.4cm]
  {\Large Sistema de Alerta Temprana y Atención al Rendimiento Académico}\\[1.2cm]
  \rule{0.6\linewidth}{1pt}\\[0.6cm]
  {\LARGE\bfseries Requisitos de Despliegue}\\[0.4cm]
  {\large Especificación técnica de infraestructura, configuración y puesta en producción}\\[0.6cm]
  \rule{0.6\linewidth}{1pt}\\[1.8cm]
  \begin{tabular}{ll}
    \textbf{Versión}     & 1.0 \\[4pt]
    \textbf{Fecha}       & Mayo 2026 \\[4pt]
    \textbf{Plataforma}  & Railway (backend) · Vercel (frontend) \\[4pt]
    \textbf{Repositorio} & \url{https://github.com/AlphaK03/ATARA-Backend} \\
  \end{tabular}
  \vfill
  {\small Ministerio de Educación Pública de Costa Rica --- División de Educación Básica}
\end{titlepage}

% ── Tabla de contenidos ───────────────────────────────────────────────────
\tableofcontents
\newpage

% ═════════════════════════════════════════════════════════════════════════
\section{Introducción}

Este documento especifica todos los requisitos técnicos necesarios para desplegar la plataforma ATARA en un entorno de producción. Está dirigido a equipos de infraestructura, DevOps y desarrolladores encargados de mantener o reproducir el entorno de ejecución del sistema.

ATARA es una REST API construida sobre Spring Boot que expone servicios de gestión académica, evaluaciones por saber, alertas tempranas y reportes pedagógicos para docentes y coordinadores del MEP. El frontend es una aplicación estática servida mediante Vercel que consume la API.

\subsection{Alcance}

El documento cubre:
\begin{itemize}
  \item Arquitectura general del sistema desplegado.
  \item Requisitos de software y versiones mínimas.
  \item Configuración de variables de entorno.
  \item Proceso de despliegue en Railway y Vercel.
  \item Configuración de la base de datos y migraciones.
  \item Seguridad, CORS y tokens JWT.
  \item Verificación de salud (\textit{health check}) del servicio.
  \item Consideraciones para el entorno local de desarrollo.
\end{itemize}

% ═════════════════════════════════════════════════════════════════════════
\section{Arquitectura General}

\subsection{Componentes del sistema}

El sistema se compone de tres capas independientes que se comunican entre sí:

\begin{center}
\begin{tabular}{L{3.2cm} L{3.5cm} L{6.5cm}}
\toprule
\textbf{Componente} & \textbf{Tecnología} & \textbf{Responsabilidad} \\
\midrule
Frontend & Vite + JavaScript & Interfaz de usuario. Servida como sitio estático en Vercel. Consume la API vía HTTPS. \\
\addlinespace
Backend (API) & Spring Boot 4.0.3 / Java 17 & Lógica de negocio, autenticación JWT, exposición de endpoints REST. Desplegado en Railway. \\
\addlinespace
Base de datos & PostgreSQL 17 & Persistencia de datos. Gestionada por el plugin de base de datos de Railway. \\
\bottomrule
\end{tabular}
\end{center}

\subsection{Flujo de comunicación}

\begin{enumerate}
  \item El navegador del usuario accede al frontend servido por Vercel a través de HTTPS.
  \item El frontend realiza llamadas REST al backend en Railway mediante la variable \texttt{VITE\_API\_URL}.
  \item El backend se autentica contra PostgreSQL usando las variables de entorno inyectadas por Railway.
  \item Flyway ejecuta las migraciones de esquema al arrancar la aplicación.
\end{enumerate}

\begin{notabox}
El frontend y el backend son completamente independientes. Pueden desplegarse o actualizarse por separado sin afectar al otro componente, siempre que el contrato de la API se mantenga.
\end{notabox}

% ═════════════════════════════════════════════════════════════════════════
\section{Requisitos de Software}

\subsection{Backend}

\begin{center}
\begin{tabular}{L{5cm} L{4cm} L{4.5cm}}
\toprule
\textbf{Componente} & \textbf{Versión mínima} & \textbf{Notas} \\
\midrule
Java (JDK)             & 17 (LTS)    & Railway provee el JDK automáticamente. \\
Spring Boot            & 4.0.3       & Incluido en el JAR empaquetado. \\
Apache Maven           & 3.9+        & Para compilación local. \\
PostgreSQL             & 15+         & Versión recomendada: 17. \\
Flyway                 & Incluido    & Gestiona el esquema automáticamente. \\
\bottomrule
\end{tabular}
\end{center}

\subsection{Frontend}

\begin{center}
\begin{tabular}{L{5cm} L{4cm} L{4.5cm}}
\toprule
\textbf{Componente} & \textbf{Versión mínima} & \textbf{Notas} \\
\midrule
Node.js    & 18 LTS  & Requerido para el build con Vite. \\
npm        & 9+      & Gestor de dependencias. \\
Vite       & 5.4+    & Herramienta de build y dev server. \\
Chart.js   & 4.5+    & Dependencia de producción. \\
\bottomrule
\end{tabular}
\end{center}

\subsection{Entorno de despliegue}

\begin{center}
\begin{tabular}{L{4cm} L{4cm} L{5.5cm}}
\toprule
\textbf{Servicio} & \textbf{Proveedor} & \textbf{Requisito} \\
\midrule
Backend hosting  & Railway  & Plan Hobby o superior. Soporte para Java/Maven detectado automáticamente. \\
Frontend hosting & Vercel   & Plan Free es suficiente. Proyecto conectado al repositorio Git. \\
Base de datos    & Railway  & Plugin PostgreSQL adjunto al proyecto de Railway. \\
\bottomrule
\end{tabular}
\end{center}

% ═════════════════════════════════════════════════════════════════════════
\section{Variables de Entorno}

Todas las variables deben configurarse en el panel de Railway (para el backend) y en el panel de Vercel (para el frontend). Ningún valor sensible debe estar codificado directamente en el código fuente ni en archivos versionados.

\subsection{Backend --- Railway}

\rowcolors{2}{grisclaro}{white}
\begin{tabularx}{\textwidth}{L{4.5cm} L{4cm} X}
\toprule
\textbf{Variable} & \textbf{Requerida} & \textbf{Descripción} \\
\midrule
\texttt{PGHOST}          & Sí (auto) & Host de PostgreSQL. Railway lo inyecta automáticamente. \\
\texttt{PGPORT}          & Sí (auto) & Puerto de PostgreSQL. Railway lo inyecta automáticamente. \\
\texttt{PGDATABASE}      & Sí (auto) & Nombre de la base de datos. \\
\texttt{PGUSER}          & Sí (auto) & Usuario de la base de datos. \\
\texttt{PGPASSWORD}      & Sí (auto) & Contraseña de la base de datos. \\
\texttt{PORT}            & Sí (auto) & Puerto HTTP del servidor. Railway lo asigna. \\
\texttt{JWT\_SECRET}     & \textbf{Sí} & Clave secreta para firmar tokens JWT. Mínimo 64 caracteres aleatorios. \\
\texttt{JWT\_EXPIRATION\_MS} & No & Expiración del access token en milisegundos. Por defecto: \texttt{86400000} (24 h). \\
\texttt{JWT\_REFRESH\_EXPIRATION\_DAYS} & No & Días de validez del refresh token. Por defecto: \texttt{7}. \\
\texttt{CORS\_ALLOWED\_ORIGINS} & \textbf{Sí} & URL exacta del frontend en Vercel. Ejemplo: \texttt{https://atara.vercel.app}. \\
\texttt{SHOW\_SQL}       & No & Mostrar SQL en logs. Establecer \texttt{false} en producción. \\
\bottomrule
\end{tabularx}
\rowcolors{0}{}{}

\vspace{6pt}

\begin{warningbox}
  El valor de \texttt{JWT\_SECRET} usado en desarrollo (\texttt{application.properties}) no debe utilizarse en producción bajo ninguna circunstancia. Genere una clave con al menos 64 caracteres usando un generador de entropía seguro (por ejemplo: \texttt{openssl rand -hex 64}).
\end{warningbox}

\subsection{Frontend --- Vercel}

\begin{center}
\begin{tabular}{L{4.5cm} L{2.5cm} L{6.5cm}}
\toprule
\textbf{Variable} & \textbf{Requerida} & \textbf{Descripción} \\
\midrule
\texttt{VITE\_API\_URL} & \textbf{Sí} & URL base de la API en Railway. Ejemplo: \texttt{https://atara-api.up.railway.app}. No incluir barra al final. \\
\bottomrule
\end{tabular}
\end{center}

\begin{notabox}
En Vercel, las variables de entorno que deben estar disponibles en el navegador deben tener el prefijo \texttt{VITE\_}. Las variables sin ese prefijo son inaccesibles desde el código JavaScript compilado.
\end{notabox}

% ═════════════════════════════════════════════════════════════════════════
\section{Base de Datos}

\subsection{Motor y versión}

El sistema utiliza \textbf{PostgreSQL 17}. En Railway se provisiona a través del plugin oficial de PostgreSQL, que crea automáticamente las credenciales y las inyecta como variables de entorno en el servicio del backend.

\subsection{Gestión del esquema con Flyway}

El esquema de la base de datos es gestionado exclusivamente por \textbf{Flyway}. Hibernate está configurado con \texttt{ddl-auto=none}, lo que significa que nunca modifica el esquema directamente.

Las migraciones se ubican en:

\begin{lstlisting}
src/main/resources/db/migration/
  V1__init_schema.sql        -- 20 tablas, 40+ índices, 13 triggers de auditoría
  V2__sample_data.sql        -- Datos semilla (credenciales de prueba)
  V3__queries_reference.sql  -- 8 vistas de reporte
  V4__evaluacion_saberes_alertas_tematicas.sql  -- 6 tablas adicionales
\end{lstlisting}

Flyway ejecuta las migraciones pendientes automáticamente al iniciar la aplicación. Si el checksum de una migración existente cambia, la aplicación \textbf{no arrancará}.

\begin{warningbox}
  Nunca modifique los archivos de migración existentes (\texttt{V1} a \texttt{V4}) una vez que hayan sido ejecutados en producción. Para cualquier cambio de esquema, cree un nuevo archivo con el siguiente número de versión (\texttt{V5\_\_descripcion.sql}).
\end{warningbox}

\subsection{Usuarios y contraseñas}

Las contraseñas de usuarios de la aplicación se almacenan hasheadas con \textbf{BCrypt} (strength 12). Los hashes del archivo \texttt{V2\_\_sample\_data.sql} son datos de prueba y deben ser regenerados para el entorno de producción.

Para generar un hash válido en Java:
\begin{lstlisting}[language=Java]
new BCryptPasswordEncoder(12).encode("contraseña_segura_aqui");
\end{lstlisting}

% ═════════════════════════════════════════════════════════════════════════
\section{Despliegue del Backend en Railway}

\subsection{Proceso de despliegue}

Railway detecta automáticamente que el proyecto es una aplicación Maven/Spring Boot y ejecuta el siguiente proceso:

\begin{enumerate}
  \item \textbf{Build:} \texttt{./mvnw clean package -DskipTests}
  \item \textbf{Inicio:} \texttt{java -jar target/atara-api-0.0.1-SNAPSHOT.jar}
  \item \textbf{Migraciones:} Flyway se ejecuta automáticamente antes de que el servidor acepte tráfico.
  \item \textbf{Health check:} Railway sondea \texttt{GET /actuator/health} hasta recibir \texttt{HTTP 200}.
\end{enumerate}

\subsection{Configuración del health check}

El archivo \texttt{railway.toml} en la raíz del repositorio define:

\begin{lstlisting}[language=bash]
[deploy]
healthcheckPath = "/actuator/health"
healthcheckTimeout = 300
\end{lstlisting}

El endpoint devuelve:
\begin{lstlisting}[language=json]
{ "status": "UP" }
\end{lstlisting}

El timeout de 300 segundos permite que Flyway complete las migraciones en bases de datos con datos existentes antes de que Railway marque el servicio como fallido.

\subsection{Puerto del servidor}

El backend lee el puerto desde la variable \texttt{PORT} inyectada por Railway. El valor por defecto para desarrollo local es \texttt{8081}:

\begin{lstlisting}
server.port=${PORT:8081}
\end{lstlisting}

No configure un puerto fijo en producción; Railway asigna el puerto dinámicamente.

% ═════════════════════════════════════════════════════════════════════════
\section{Despliegue del Frontend en Vercel}

\subsection{Configuración del proyecto}

El frontend reside en el subdirectorio \texttt{temp-frontend/} del repositorio. En Vercel debe configurarse:

\begin{center}
\begin{tabular}{L{4.5cm} L{8cm}}
\toprule
\textbf{Parámetro Vercel} & \textbf{Valor} \\
\midrule
Root Directory      & \texttt{temp-frontend} \\
Build Command       & \texttt{npm run build} \\
Output Directory    & \texttt{dist} \\
Install Command     & \texttt{npm install} \\
Node.js Version     & 18.x \\
\bottomrule
\end{tabular}
\end{center}

\subsection{Proceso de despliegue}

Vercel está conectado directamente al repositorio de GitHub. Cada \textit{push} a la rama \texttt{main} dispara un nuevo despliegue automático. El proceso es:

\begin{enumerate}
  \item \texttt{npm install} instala las dependencias.
  \item \texttt{vite build} compila y empaqueta los archivos estáticos en \texttt{dist/}.
  \item Vercel sirve el contenido de \texttt{dist/} como una CDN global.
\end{enumerate}

\begin{notabox}
  La variable \texttt{VITE\_API\_URL} debe estar definida en Vercel \textbf{antes} del build. Vite incrusta las variables \texttt{VITE\_*} en el bundle en tiempo de compilación, no en tiempo de ejecución. Si cambia la URL del backend, debe re-desplegar el frontend.
\end{notabox}

% ═════════════════════════════════════════════════════════════════════════
\section{Seguridad}

\subsection{Autenticación JWT}

El backend implementa autenticación basada en JSON Web Tokens. El flujo es:

\begin{enumerate}
  \item El cliente realiza \texttt{POST /api/auth/login} con correo y contraseña.
  \item El servidor devuelve un \textbf{access token} (duración configurable, por defecto 24 h) y un \textbf{refresh token} (por defecto 7 días).
  \item El cliente incluye el access token en el header \texttt{Authorization: Bearer <token>} en cada solicitud protegida.
  \item Al expirar el access token, el cliente renueva usando \texttt{POST /api/auth/refresh}.
\end{enumerate}

Los únicos endpoints públicos son:
\begin{itemize}
  \item \texttt{POST /api/auth/login}
  \item \texttt{POST /api/auth/refresh}
  \item \texttt{GET /actuator/health}
\end{itemize}

\subsection{CORS}

La variable \texttt{CORS\_ALLOWED\_ORIGINS} debe contener la URL exacta del frontend en producción. Usar el valor \texttt{*} en producción está \textbf{prohibido} ya que permitiría solicitudes desde cualquier origen.

Ejemplo de configuración correcta:
\begin{lstlisting}
CORS_ALLOWED_ORIGINS=https://atara.vercel.app
\end{lstlisting}

\subsection{Buenas prácticas de seguridad}

\begin{itemize}
  \item Nunca exponer \texttt{spring.jpa.show-sql=true} en producción: establezca \texttt{SHOW\_SQL=false}.
  \item El endpoint \texttt{/actuator/health} debe mostrar únicamente el estado (\texttt{show-details=never}).
  \item No registrar en logs contraseñas, tokens, cédulas ni nombres completos combinados con calificaciones.
  \item Las contraseñas deben compararse únicamente mediante \texttt{BCryptPasswordEncoder}; nunca en texto plano.
  \item Utilizar HTTPS en todos los endpoints de producción (Railway y Vercel lo proveen por defecto).
\end{itemize}

% ═════════════════════════════════════════════════════════════════════════
\section{Entorno de Desarrollo Local}

\subsection{Requisitos previos}

\begin{itemize}
  \item JDK 17 instalado y \texttt{JAVA\_HOME} configurado.
  \item Maven 3.9+ (o usar el wrapper incluido: \texttt{./mvnw}).
  \item Node.js 18 LTS y npm 9+.
  \item Docker Desktop (para levantar PostgreSQL en contenedor).
\end{itemize}

\subsection{Pasos para iniciar el entorno}

\begin{enumerate}
  \item \textbf{Levantar la base de datos:}
\begin{lstlisting}[language=bash]
docker-compose up -d
\end{lstlisting}
  Esto inicia un contenedor PostgreSQL 17 con las siguientes credenciales por defecto:
  \begin{itemize}
    \item Base de datos: \texttt{atara\_db}
    \item Usuario: \texttt{atara\_user}
    \item Contraseña: \texttt{atara\_pass\_123}
    \item Puerto local: \texttt{5433} (mapeado al puerto interno 5432)
  \end{itemize}

  \item \textbf{Iniciar el backend:}
\begin{lstlisting}[language=bash]
./mvnw spring-boot:run
\end{lstlisting}
  El servidor arranca en \texttt{http://localhost:8081}. Flyway ejecuta las migraciones automáticamente en el primer inicio.

  \item \textbf{Iniciar el frontend:}
\begin{lstlisting}[language=bash]
cd temp-frontend
npm install
npm run dev
\end{lstlisting}
  La aplicación queda disponible en \texttt{http://localhost:3000}. El servidor de desarrollo de Vite proxifica las llamadas a \texttt{/api} y \texttt{/actuator} hacia \texttt{http://localhost:8081}.
\end{enumerate}

\subsection{Diferencias respecto a producción}

\begin{center}
\begin{tabular}{L{4.5cm} L{3.5cm} L{5cm}}
\toprule
\textbf{Aspecto} & \textbf{Local} & \textbf{Producción} \\
\midrule
Puerto backend     & 8081         & Asignado por Railway \\
Puerto PostgreSQL  & 5433         & Interno en Railway \\
Puerto frontend    & 3000         & CDN global (Vercel) \\
Logs SQL           & Activados    & Desactivados \\
CORS               & \texttt{*}   & Dominio Vercel específico \\
JWT secret         & Valor de prueba & Clave segura de 64+ chars \\
HTTPS              & No           & Sí (obligatorio) \\
\bottomrule
\end{tabular}
\end{center}

% ═════════════════════════════════════════════════════════════════════════
\section{Verificación Post-Despliegue}

Una vez completado el despliegue, verificar los siguientes puntos:

\begin{enumerate}
  \item \textbf{Health check del backend:}
\begin{lstlisting}[language=bash]
curl https://<dominio-railway>/actuator/health
# Respuesta esperada: {"status":"UP"}
\end{lstlisting}

  \item \textbf{Login desde el frontend:} Acceder a la URL de Vercel y verificar que el formulario de inicio de sesión se comunica correctamente con el backend.

  \item \textbf{Migraciones:} Confirmar en los logs de Railway que Flyway reporta todas las migraciones como \texttt{Success} sin errores de checksum.

  \item \textbf{CORS:} Abrir las herramientas de desarrollador del navegador y verificar que no aparecen errores de CORS en la consola al hacer login o navegar.

  \item \textbf{Variables de entorno:} En Railway, confirmar que las variables \texttt{PGHOST}, \texttt{PGPORT}, \texttt{PGDATABASE}, \texttt{PGUSER}, \texttt{PGPASSWORD} y \texttt{JWT\_SECRET} están presentes antes del primer despliegue.
\end{enumerate}

% ═════════════════════════════════════════════════════════════════════════
\section{Referencia de Puertos y URLs}

\begin{center}
\begin{tabular}{L{5cm} L{3cm} L{5.5cm}}
\toprule
\textbf{Servicio} & \textbf{Puerto/URL} & \textbf{Contexto} \\
\midrule
Backend (local)         & \texttt{:8081}  & Desarrollo local \\
Frontend dev (local)    & \texttt{:3000}  & Desarrollo local \\
PostgreSQL (local)      & \texttt{:5433}  & Mapeado desde Docker \\
Backend (Railway)       & Dinámico (HTTPS) & Producción \\
Frontend (Vercel)       & HTTPS           & Producción \\
Health check            & \texttt{/actuator/health} & Ambos entornos \\
\bottomrule
\end{tabular}
\end{center}

% ═════════════════════════════════════════════════════════════════════════
\section{Dependencias Clave del Backend}

\begin{center}
\begin{tabular}{L{5.5cm} L{2.5cm} L{5.5cm}}
\toprule
\textbf{Dependencia} & \textbf{Versión} & \textbf{Propósito} \\
\midrule
spring-boot-starter-web       & 4.0.3  & API REST \\
spring-boot-starter-security  & 4.0.3  & Seguridad y autenticación \\
spring-boot-starter-data-jpa  & 4.0.3  & Acceso a datos con Hibernate \\
spring-boot-starter-validation & 4.0.3 & Validación de DTOs \\
spring-boot-starter-actuator  & 4.0.3  & Monitoreo y health check \\
spring-boot-starter-flyway    & 4.0.3  & Migraciones de base de datos \\
postgresql (JDBC driver)      & Incluido & Conexión a PostgreSQL \\
jjwt-api / jjwt-impl          & 0.12.6 & Generación y validación de JWT \\
lombok                        & Incluido & Reducción de código repetitivo \\
pdfbox                        & 3.0.3  & Procesamiento PDF (importar PIAD) \\
tess4j                        & 5.11.0 & OCR para extracción de listas PIAD \\
\bottomrule
\end{tabular}
\end{center}

\vspace{1cm}
\hrule
\vspace{0.4cm}
{\small\textit{Documento generado para el proyecto ATARA --- MEP Costa Rica. Para cambios en la infraestructura o requisitos adicionales, actualizar este documento junto con el código correspondiente.}}

\end{document}
